\chapter{La réaction de Wittig et les autres formations de C=C}\label{ch:b2:wittig}

La mouche domestique femelle attire les mâles par un hydrocarbure à 23 carbones porteur d'une seule
double liaison, entre les carbones 9 et 10, de géométrie (\textit{Z}). Pour le produire en quantité,
un chimiste doit placer cette unique double liaison exactement à cet endroit, et avec cette
géométrie~: une élimination la disperserait sur plusieurs positions et donnerait les deux géométries.
Ce chapitre construit les doubles liaisons par l'autre voie, en réunissant deux fragments au niveau
d'un groupe carbonyle, de sorte que la double liaison apparaisse là où se trouvait le \ce{C=O}.
% ledger: use:muscalure

\begin{recall}
\Cref{ch:b2:enolates-aldol}~: la \omterm{def:b2:enolates-aldol:aldol}{crotonisation}. \Cref{ch:b2:alkene-redox}~: l'hydrogénation d'un
alcyne en alcène (\textit{Z}) sur un catalyseur empoisonné. Le volume de première année~: E1, E2 et
règle de Zaïtsev, substitution bimoléculaire, organomagnésiens et inversion de polarité, descripteurs
(\textit{Z})/(\textit{E}). Le volume scolaire~: économie d'atomes.
\end{recall}

\section{Sels de phosphonium et ylures}

\begin{definition}[Ylures de phosphonium]\label{def:b2:wittig:ylide}
Un \emph{sel de phosphonium}\index{sel de phosphonium} $\mathrm{R_3P^+{-}CH_2R'\ X^-}$ se prépare par
substitution bimoléculaire d'un halogénoalcane par une phosphine, le plus souvent la triphénylphosphine
\ce{Ph3P}. Arracher un hydrogène au carbone lié au phosphore donne un \emph{ylure de phosphonium}\index{ylure de phosphonium},
molécule neutre portant deux charges opposées voisines, $\mathrm{Ph_3P^+{-}CH^-{-}R'}$, qu'on écrit aussi
$\mathrm{Ph_3P{=}CHR'}$. Un \emph{ylure stabilisé}\index{ylure stabilise@ylure stabilisé} porte sur ce
carbone un groupe attracteur d'électrons (un ester, une cétone, un \omterm{def:b2:amines:nitrile}{nitrile}) qui délocalise la charge
négative.
\end{definition}

\begin{omfigure}
\setchemfig{atom sep=1.9em}
\schemestart
\chemfig{Ph_3P} \+ \chemfig{CH_3-I}
\arrow{->[S$_\mathrm{N}$2]}[,0.9]
\chemfig{Ph_3P^{+}-CH_3}\;\chemfig{I^{-}}
\arrow{->[\ce{BuLi}][$-$ \ce{BuH}, \ce{LiI}]}[,1.3]
\chemfig{Ph_3P^{+}-\chemabove{C}{\scriptstyle -}H_2}
\arrow{<->}[,0.8]
\chemfig{Ph_3P=CH_2}
\schemestop
\par\vspace{6pt}

{\small Le méthylidènetriphénylphosphorane, l'ylure le plus simple. La triphénylphosphine déplace
l'iodure de l'iodométhane~; le butyllithium arrache un proton au groupe méthyle. Les deux structures de
droite sont deux écritures de la même molécule~: celle à charges séparées est la plus proche de la réalité.}
\end{omfigure}

\begin{proposition}[Préparer l'ylure]\label{prop:b2:wittig:ylide-acidity}
Un \omterm{def:b2:wittig:ylide}{sel de phosphonium} est bien plus acide sur le carbone voisin du phosphore qu'un alcane. Un sel
d'alkylphosphonium simple exige néanmoins une base très forte (butyllithium, hydrure de sodium,
amidure de sodium), en conditions anhydres et inertes~; un sel dont le \ce{CH2} porte aussi un
groupe carbonyle est déprotoné par une base faible (carbonate de sodium, hydroxyde de sodium), et
l'\omterm{def:b2:wittig:ylide}{ylure stabilisé} formé peut être isolé et conservé.
\end{proposition}

\begin{proof}[Argument]
Le phosphore positif voisin du carbanion le stabilise par sa charge et par la polarisabilité de ce
gros atome~: un ylure simple est un carbanion stabilisé, mais moins qu'un \omterm{def:b2:enolates-aldol:enolate}{énolate}, d'où la base forte.
Un carbonyle voisin répartit la charge négative sur l'oxygène, comme dans un \omterm{def:b2:enolates-aldol:enolate}{énolate}
(\cref{prop:b2:enolates-aldol:alpha-acidity}), en plus du phosphore~: l'acide conjugué devient assez
acide pour une base faible, et l'ylure, bien moins basique, ne réagit plus avec l'eau ni avec l'air.
\end{proof}

\section{La réaction de Wittig}

\begin{definition}[Réaction de Wittig]\label{def:b2:wittig:wittig}
La \emph{réaction de Wittig}\index{reaction de Wittig@réaction de Wittig} d'un \omterm{def:b2:wittig:ylide}{ylure de phosphonium}
avec un aldéhyde ou une cétone donne un alcène et l'oxyde de triphénylphosphine~:
\begin{equation*}
\mathrm{Ph_3P{=}CHR + R'CHO \longrightarrow R'CH{=}CHR + Ph_3P{=}O}.
\end{equation*}
Elle passe par un \emph{oxaphosphétane}\index{oxaphosphetane@oxaphosphétane}, cycle à quatre atomes
formé du phosphore, de deux carbones et de l'oxygène, qui apparaît quand le carbone de l'ylure se lie
au carbone du carbonyle tandis que l'oxygène du carbonyle se lie au phosphore.
\end{definition}

\begin{omfigure}
\setchemfig{atom sep=1.9em}
\schemestart
\chemfig{Ph_3P=CH-R} \+ \chemfig{O=CH-R'}
\arrow{->[addition]}[,1.0]
\chemfig{Ph_3P?-[6]CH(-[4]R)-[0]CH(-[0]R')-[2]O?}
\schemestop

\medskip
\schemestart
\arrow{->[élimination]}[,1.1]
\chemfig{R-CH=CH-R'} \+ \chemfig{Ph_3P=O}
\schemestop
\par\vspace{6pt}

{\small La \omterm{def:b2:wittig:wittig}{réaction de Wittig}. Le carbone de l'ylure et le carbone du carbonyle se lient tandis que
l'oxygène se lie au phosphore, en une seule étape~: c'est l'\omterm{def:b2:wittig:wittig}{oxaphosphétane}. Son cycle s'ouvre ensuite
dans l'autre sens, en rompant les liaisons \ce{P-C} et \ce{C-O}~: la liaison \ce{C=C} se forme entre
les deux anciens carbones réactifs, et la liaison \ce{P=O} de l'oxyde de triphénylphosphine se forme
en même temps.}
\end{omfigure}

\begin{proposition}[Une double liaison exactement à sa place]\label{prop:b2:wittig:regiospecific}
Dans une \omterm{def:b2:wittig:wittig}{réaction de Wittig}, la nouvelle liaison \ce{C=C} relie l'ancien carbone du carbonyle à
l'ancien carbone de l'ylure, et nulle part ailleurs.
\end{proposition}

\begin{proof}[Argument]
Les seules liaisons formées et rompues sont celles du cycle à quatre atomes~: le carbone de l'ylure
avec le carbone du carbonyle, puis les liaisons \ce{C-P} et \ce{C-O}. Aucun hydrogène ne se déplace et
aucun carbocation ni carbanion ne se forme le long de la chaîne, si bien que la double liaison ne peut
pas migrer. Une élimination, au contraire, choisit entre les hydrogènes de plusieurs carbones voisins
(règle de Zaïtsev), et une réaction E1 peut s'accompagner d'un réarrangement.
\end{proof}

\begin{proposition}[Force motrice]\label{prop:b2:wittig:driving}
La formation de la liaison \ce{P=O} de l'oxyde de triphénylphosphine rend la \omterm{def:b2:wittig:wittig}{réaction de Wittig}
fortement favorable et irréversible.
\end{proposition}

\begin{proof}[Argument]
La réaction échange une liaison \ce{C=O} et une liaison \ce{P-C} (dans l'ylure, la \ce{P=C} est
surtout une liaison simple \ce{P+-C-}) contre une liaison \ce{C=C} et une liaison \ce{P=O}. La liaison
\ce{P=O} d'un oxyde de phosphine est très forte, bien plus forte que la liaison \ce{P-C} perdue, et la
perte du \ce{C=O} est compensée en partie par la formation du \ce{C=C}~: le bilan est nettement
favorable. Le phosphore est le puits à oxygène de la réaction.
\end{proof}

\section{Stéréosélectivité}

\begin{proposition}[Géométrie de l'alcène]\label{prop:b2:wittig:stereo}
Avec un aldéhyde, un ylure non stabilisé (\ce{Ph3P=CHR}, R un groupe alkyle) donne surtout l'alcène
(\textit{Z}), en l'absence de sels. Un \omterm{def:b2:wittig:ylide}{ylure stabilisé} (\ce{Ph3P=CH-COOR}, \ce{Ph3P=CH-COR}) donne
surtout l'alcène (\textit{E}). Un \omterm{def:b2:wittig:ylide}{ylure stabilisé} seulement par un groupe aryle donne des mélanges.
\end{proposition}

\admitted

Ce sont des observations. Leur explication habituelle lit la \omterm{def:b2:wittig:wittig}{réaction de Wittig} comme une compétition
entre vitesses et stabilités. Avec un ylure réactif, non stabilisé, l'\omterm{def:b2:wittig:wittig}{oxaphosphétane} se forme vite et
de façon irréversible, par l'état de transition le moins encombré, dans lequel les deux substituants se
retrouvent du même côté du cycle~; il se fragmente ensuite en conservant cette disposition~: le contrôle
cinétique donne l'alcène (\textit{Z}). Avec un \omterm{def:b2:wittig:ylide}{ylure stabilisé}, la première étape est plus lente et
renversable~; les \omterm{def:b2:wittig:wittig}{oxaphosphétanes} s'interconvertissent en repassant par les réactifs, et celui dont les
substituants sont de part et d'autre, moins tendu, l'emporte~: l'alcène (\textit{E}), sous contrôle
thermodynamique.

\begin{omfigure}
\setchemfig{atom sep=1.9em}
\schemestart
\chemfig{Ph_3P=CH-CH_2CH_3} \+ \chemfig{O=CH-Ph}
\arrow{->}[,0.8]
\chemfig{Ph-[:-30]=[:30]-[:90]CH_2CH_3}
\schemestop

\medskip
\schemestart
\chemfig{Ph_3P=CH-COOEt} \+ \chemfig{O=CH-Ph}
\arrow{->}[,0.8]
\chemfig{Ph-[:-30]=[:30]-[:-30]COOEt}
\schemestop
\par\vspace{6pt}

{\small Deux ylures, un aldéhyde. En haut~: l'ylure de propylidène, non stabilisé, donne surtout le
(\textit{Z})-1-phénylbut-1-ène. En bas~: l'\omterm{def:b2:wittig:ylide}{ylure stabilisé} donne surtout le
(\textit{E})-3-phénylpropénoate d'éthyle (cinnamate d'éthyle). L'oxyde de triphénylphosphine, formé
dans les deux cas, n'est pas représenté.}
\end{omfigure}

\section{La réaction de Horner--Wadsworth--Emmons}

\begin{definition}[Phosphonates et réaction HWE]\label{def:b2:wittig:hwe}
Un \emph{phosphonate}\index{phosphonate} est un ester d'un acide phosphonique, $\mathrm{(RO)_2P({=}O){-}R'}$,
qui comporte une liaison \ce{P-C}. Il se prépare par la réaction de Michaelis--Arbuzov d'un phosphite de
trialkyle avec un halogénoalcane, $\mathrm{P(OEt)_3 + R'Br \to (EtO)_2P(O)R' + EtBr}$. Dans la
\emph{réaction de Horner--Wadsworth--Emmons}\index{reaction de Horner--Wadsworth--Emmons@réaction de Horner--Wadsworth--Emmons} (HWE), l'anion
d'un phosphonate portant un groupe attracteur d'électrons sur son \ce{CH2} s'additionne sur un aldéhyde
et donne un alcène et un anion phosphate de dialkyle.
\end{definition}

\begin{omfigure}
\setchemfig{atom sep=1.8em}
\schemestart
\chemfig{EtO-P(=[2]O)(-[6]OEt)-CH_2-COOEt}
\arrow{->[\ce{NaH}][$-$ \ce{H2}]}[,0.9]
\chemfig{EtO-P(=[2]O)(-[6]OEt)-\chemabove{C}{\scriptstyle -}H-COOEt}
\schemestop

\medskip
\schemestart
\arrow{->[\ce{PhCHO}]}[,0.9]
\chemfig{*6(=-=(-[:30]=[:-30]-[:30](=[:90]O)-[:-30]O-[:30]Et)-=-)}
\+
\chemfig{EtO-P(=[2]O)(-[6]O^{-})-OEt}
\schemestop
\par\vspace{6pt}

{\small La \omterm{def:b2:wittig:hwe}{réaction de Horner--Wadsworth--Emmons} du phosphonoacétate de triéthyle avec le benzaldéhyde.
L'hydrure de sodium arrache l'hydrogène situé entre le \omterm{def:b2:wittig:hwe}{phosphonate} et l'ester~; l'anion s'additionne sur
l'aldéhyde, et l'intermédiaire à quatre atomes se fragmente comme dans la \omterm{def:b2:wittig:wittig}{réaction de Wittig}. Le produit
est le (\textit{E})-cinnamate d'éthyle~; le sous-produit est l'anion phosphate de diéthyle, soluble dans l'eau.}
\end{omfigure}

\begin{proposition}[Pourquoi les chimistes aiment la réaction HWE]\label{prop:b2:wittig:hwe}
La réaction HWE donne surtout l'alcène (\textit{E}), et son sous-produit phosphoré, un sel soluble dans
l'eau, s'élimine par un simple lavage aqueux~; l'anion \omterm{def:b2:wittig:hwe}{phosphonate} est en outre plus nucléophile que
l'\omterm{def:b2:wittig:ylide}{ylure stabilisé} correspondant et réagit aussi avec les cétones.
\end{proposition}

\begin{proof}[Argument]
L'anion \omterm{def:b2:wittig:hwe}{phosphonate} est un carbanion stabilisé, si bien que son addition est renversable et que la
réaction se déroule sous contrôle thermodynamique~: (\textit{E}), comme avec les \omterm{def:b2:wittig:ylide}{ylures stabilisés}.
Contrairement à un \omterm{def:b2:wittig:ylide}{ylure de phosphonium}, il porte une charge négative entière, que ne compense aucun
cation voisin~: il est plus réactif. Le sous-produit, \ce{(EtO)2PO2-}, est un sel ionique, alors que
l'oxyde de triphénylphosphine est un solide organique neutre qu'il faut séparer du produit par
chromatographie ou par cristallisation.
\end{proof}

\section{Choisir une méthode}

\begin{method}[Déconnexion de Wittig]\label{met:b2:wittig:wittig-disconnection}
\begin{enumerate}
  \item Couper la liaison \ce{C=C} cible~; l'un des carbones devient un carbone de carbonyle, l'autre le
        carbone de l'ylure (un \omterm{def:b2:wittig:ylide}{sel de phosphonium}, issu d'un halogénure).
  \item Des deux façons de procéder, préférer celle dont l'halogénure est primaire (substitution
        bimoléculaire par \ce{Ph3P}) et dont le composé carbonylé est un aldéhyde.
  \item Choisir le réactif selon la géométrie voulue~: un ylure non stabilisé pour (\textit{Z}), un
        \omterm{def:b2:wittig:ylide}{ylure stabilisé} ou un \omterm{def:b2:wittig:hwe}{phosphonate} (HWE) pour (\textit{E}).
  \item Vérifier que rien d'autre dans la molécule ne réagit avec la base utilisée.
\end{enumerate}
\end{method}

\begin{method}[Choisir une réaction formant une liaison C=C]\label{met:b2:wittig:choose-cc}
\begin{center}
\small
\begin{tabular}{@{}lll@{}}
\hline
Méthode & Position de la \ce{C=C} & Géométrie \\
\hline
Élimination E2 ou E1 & règle de Zaïtsev, mélanges & surtout (\textit{E}), mélanges \\
\omterm{def:b2:enolates-aldol:aldol}{Crotonisation} & conjuguée avec \ce{C=O} & (\textit{E}) \\
Wittig, ylure non stabilisé & à la place de l'ancien \ce{C=O} & (\textit{Z}) \\
Wittig, \omterm{def:b2:wittig:ylide}{ylure stabilisé}~; HWE & à la place de l'ancien \ce{C=O} & (\textit{E}) \\
Alcyne, puis catalyseur empoisonné & à la place de l'ancienne triple liaison & (\textit{Z}) \\
\hline
\end{tabular}
\end{center}
\medskip
Pour réunir deux fragments de taille comparable au niveau d'une \ce{C=C}, la \omterm{def:b2:wittig:wittig}{réaction de Wittig} et
ses variantes sont en général le premier choix~; la métathèse des oléfines, qui échange les extrémités
de deux alcènes, est traitée dans le volume de troisième année.
\end{method}

\begin{history}[L'ylure devenu réactif]
Georg Wittig, qui étudiait les ylures du phosphore au début des années 1950, découvrit que le
méthylidènetriphénylphosphorane transforme la benzophénone en 1,1-diphényléthène et en oxyde de
triphénylphosphine~; avec son étudiant Ulrich Schöllkopf, il fit de cette observation une méthode
générale en 1954. L'industrie l'adopta rapidement, notamment pour la vitamine A, et Wittig reçut le
prix Nobel de chimie 1979, partagé.
\end{history}

\begin{safety}{\ghs[5mm]{GHS02}\,\ghs[5mm]{GHS05}\,\ghs[5mm]{GHS07}\,\ghs[5mm]{GHS08}}
Les solutions de butyllithium sont inflammables et corrosives et réagissent violemment avec l'eau~;
l'hydrure de sodium libère du dihydrogène au contact de l'eau. La triphénylphosphine est nocive et
corrosive pour les yeux, le phosphite de triéthyle inflammable et nocif~; l'oxyde de
triphénylphosphine est nocif.
% ledger: ghs:BuLi, ghs:PPh3, ghs:triethyl-phosphite, ghs:Ph3PO
\end{safety}

\section{Exercices}

\begin{exercise}[$\star$]\label{exo:b2:wittig:1}
Écrire la préparation de l'ylure \ce{Ph3P=CHCH3} à partir de la triphénylphosphine et d'un halogénoalcane.
\end{exercise}

\begin{exercise}[$\star$]\label{exo:b2:wittig:2}
Donner les produits de la \omterm{def:b2:wittig:wittig}{réaction de Wittig} du benzaldéhyde avec l'ylure préparé à partir de l'iodométhane.
\end{exercise}

\begin{exercise}[$\star$]\label{exo:b2:wittig:3}
Classer en stabilisé, non stabilisé ou stabilisé par un aryle~: \ce{Ph3P=CHCOOEt}, \ce{Ph3P=CHCH3},
\ce{Ph3P=CHCOCH3}, \ce{Ph3P=CHC6H5}. Lequel exige le butyllithium~?
\end{exercise}

\begin{exercise}[$\star$]\label{exo:b2:wittig:4}
Donner le produit du phosphonoacétate de triéthyle, de l'hydrure de sodium et du propanal, avec sa géométrie.
\end{exercise}

\begin{exercise}[$\star\star$]\label{exo:b2:wittig:5}
On veut préparer le méthylènecyclohexane. Comparer une voie de Wittig à partir de la cyclohexanone et
la déshydratation en milieu acide du 1-méthylcyclohexan-1-ol.
\end{exercise}

\begin{exercise}[$\star\star$]\label{exo:b2:wittig:6}
Déconnecter le (\textit{E})-1,2-diphényléthène (stilbène) par la \omterm{def:b2:wittig:wittig}{réaction de Wittig}. Pourquoi la voie
de Wittig est-elle un mauvais choix pour le (\textit{E})-stilbène pur, et que faudrait-il utiliser à la place~?
\end{exercise}

\begin{exercise}[$\star\star$]\label{exo:b2:wittig:7}
Le propanal réagit d'une part avec \ce{Ph3P=CHCH3}, d'autre part avec \ce{Ph3P=CHCOOEt}. Donner les
produits et leur géométrie principale.
\end{exercise}

\begin{exercise}[$\star\star$]\label{exo:b2:wittig:8}
Calculer l'économie d'atomes de la méthylénation de la cyclohexanone par \ce{Ph3P=CH2}.
\end{exercise}

\begin{exercise}[$\star\star$]\label{exo:b2:wittig:9}
Écrire la réaction de Michaelis--Arbuzov du phosphite de triéthyle avec le bromoéthanoate d'éthyle, en
deux étapes~: substitution par le phosphore, puis désalkylation par le bromure.
\end{exercise}

\begin{exercise}[$\star\star\star$]\label{exo:b2:wittig:10}
Proposer une synthèse du 1,6-diphénylhexa-1,3,5-triène à partir du (\textit{E})-but-2-ènedial
\ce{OHC-CH=CH-CHO} par deux \omterm{def:b2:wittig:wittig}{réactions de Wittig}, en précisant l'ylure et la quantité nécessaire.
\end{exercise}

\begin{exercise}[$\star\star\star$]\label{exo:b2:wittig:11}
La (\textit{E})-4-phénylbut-3-én-2-one peut être préparée par une \omterm{def:b2:enolates-aldol:aldol}{crotonisation} ou par une \omterm{def:b2:wittig:wittig}{réaction
de Wittig}. Écrire les deux voies et les comparer.
\end{exercise}

\begin{exercise}[$\star\star\star$]\label{exo:b2:wittig:12}
Proposer deux voies vers le (\textit{Z})-oct-4-ène, l'une par la \omterm{def:b2:wittig:wittig}{réaction de Wittig} et l'autre par un
alcyne, et dire quels sous-produits laisse chacune.
\end{exercise}

\section{Problème~: une phéromone par la réaction de Wittig}

\begin{problem}[{Problème du week-end --- la rétrosynthèse d'un alcène (Z), l'ylure non stabilisé et
sa sélectivité, le coût en oxyde de triphénylphosphine, et la variante par un alcyne}]
\label{pb:b2:wittig:1}
Cible~: le (\textit{Z})-tricos-9-ène, \ce{CH3(CH2)7CH=CH(CH2)12CH3}, la phéromone de la mouche
domestique de l'introduction. Masses molaires (\unit{g/mol})~: C 12.011, H 1.008, O 15.999, P 30.974, Br 79.904.
% ledger: use:muscalure, aw:C, aw:H, aw:O, aw:P, aw:Br, ghs:nonanal, ghs:BuLi

\medskip
\noindent\textbf{Partie I --- Rétrosynthèse.}
\begin{enumerate}
  \item Écrire la formule brute de la cible et vérifier qu'il s'agit d'un alcène.
  \item Quelle liaison une déconnexion de Wittig coupe-t-elle~?
  \item Donner les deux couples possibles (composé carbonylé, \omterm{def:b2:wittig:ylide}{sel de phosphonium}).
  \item Pour le couple nonanal + ylure, quel halogénoalcane fournit le \omterm{def:b2:wittig:ylide}{sel de phosphonium}~?
  \item Pourquoi un halogénoalcane primaire convient-il à cette étape~?
  \item Écrire la formation du \omterm{def:b2:wittig:ylide}{sel de phosphonium}.
\end{enumerate}

\medskip
\noindent\textbf{Partie II --- L'ylure et la géométrie.}
\begin{enumerate}[resume]
  \item Quelle base déprotone ce sel~? Pourquoi l'hydroxyde de sodium ne conviendrait-il pas~?
  \item L'ylure est-il stabilisé~? Quelle géométrie cela prévoit-il~?
  \item Écrire l'\omterm{def:b2:wittig:wittig}{oxaphosphétane} formé avec le nonanal.
  \item Pourquoi faut-il conduire la réaction à sec et sous gaz inerte~?
  \item Pourquoi la double liaison se trouve-t-elle exactement entre les carbones 9 et 10~?
  \item Que donnerait au contraire la déshydratation en milieu acide du tricosan-9-ol~?
\end{enumerate}

\medskip
\noindent\textbf{Partie III --- L'oxyde de triphénylphosphine.}
\begin{enumerate}[resume]
  \item Écrire l'équation globale à partir de l'ylure et du nonanal.
  \item Calculer les masses molaires de la cible, du nonanal, du \omterm{def:b2:wittig:ylide}{sel de phosphonium} et de l'oxyde de
        triphénylphosphine.
  \item Calculer l'économie d'atomes de l'étape de Wittig (ylure + aldéhyde).
  \item Comment sépare-t-on l'oxyde de triphénylphosphine d'un hydrocarbure~?
  \item Pourquoi la formation de l'oxyde de triphénylphosphine est-elle la force motrice~?
  \item Pourquoi la réaction HWE ne conviendrait-elle pas à cette cible~?
\end{enumerate}

\medskip
\noindent\textbf{Partie IV --- La voie par un alcyne.}
\begin{enumerate}[resume]
  \item Quel alcyne, hydrogéné, donne la cible, et sur quel catalyseur~?
  \item Construire cet alcyne à partir du déc-1-yne et d'un halogénoalcane~: quel halogénure, quelle base~?
  \item Pourquoi l'hydrogénation doit-elle s'arrêter à l'alcène, et comment le catalyseur l'assure-t-il~?
  \item Comparer les deux voies~: nombre d'étapes, maîtrise de la géométrie, sous-produits.
  \item Quelle est l'économie d'atomes de l'étape d'hydrogénation~?
  \item Donner la masse d'oxyde de triphénylphosphine coproduite par gramme de phéromone par la voie de
        Wittig.
\end{enumerate}
\end{problem}
